\chapter{Wireless Network Security}
\label{ch:3.2}
\chaptertopic{Topic 3.2}

\chapterquote{Be conservative in what you do, be liberal in what you accept from others.}{Jon Postel, RFC 761 (1980)}

\section{Motivation}

Chapter~\ref{ch:3.1} treated the network as a wired abstraction. Radio changes one thing
fundamentally: every frame a Wi-Fi device sends radiates in all directions, and any device within
range receives the same signal, whether it was invited to or not. A wireless protocol has to do the
authentication and encryption work that a wired cable gets almost for free, simply by being
physically hard to tap. This chapter covers how that protocol evolved, then steps back to the
policies that govern how people -- not just protocols -- are expected to behave. It closes with the
model that ties every control from this unit together into one coherent defense.

\begin{apalignment}{Unit 3: Securing Networks}{Topic 3.2 -- Protecting Networks: Managerial Controls and Wireless Security}
\begin{apobjectives}
\apobjective{3.2.A}{Mitigate Risk}{Identify the managerial controls that govern how a network is
secured -- written policy, standards and baselines, change control, vendor and acceptable-use
agreements, security training, and audit. Explain how each shapes what the technical controls
are actually configured to do.}
\apobjective{3.2.B}{Mitigate Risk}{Configure the security features of a wireless network
deliberately: choose between WPA2 and WPA3 and know why, require Protected Management Frames,
disable convenience features that weaken the network, and decide between a shared key and per-user
enterprise authentication with 802.1X and EAP.}
\end{apobjectives}
\end{apalignment}

\section{Wi-Fi Security Protocols}

\begin{definition}{WEP (Wired Equivalent Privacy)}
The original 1997 Wi-Fi security protocol, using RC4 with a static 40- or 104-bit key. Broken in
2001, when Fluhrer, Mantin, and Shamir showed its key-scheduling algorithm leaked enough
information to recover the key from captured traffic.\srccite{3.2}{1} A follow-up attack published
in 2007 turned that theoretical weakness into a practical one, recovering a 104-bit WEP key from as
few as 40{,}000 captured frames.\srccite{3.2}{2}
\end{definition}

WEP's failure motivated everything that followed. \textbf{WPA2}, ratified in 2004, replaced WEP's
broken cipher with AES in Counter Mode with CBC-MAC (\textbf{CCMP}), providing both confidentiality
and integrity. An attacker who captures every frame still cannot decrypt the payload or forge new
packets. WPA2 runs in two modes. \textbf{Personal} (WPA2-PSK) distributes a single pre-shared key
to everyone; each device still derives its own unique session key during a four-way handshake, so
capturing one device's traffic does not expose another's. \textbf{Enterprise} (WPA2-EAP) has each
user authenticate individually through an 802.1X server, typically RADIUS, giving every account its
own identity -- so revoking one does not affect anyone else.

WPA2's main weakness lives in that four-way handshake: an attacker who captures it can take it
offline and run a dictionary or brute-force attack against the PSK. There is no rate limit at all,
because the attack happens entirely on the attacker's own hardware. A weak or guessable password
is recovered this way regardless of how strong AES-CCMP itself is.

\begin{definition}{SAE (Simultaneous Authentication of Equals)}
WPA3's replacement handshake, also called Dragonfly: a zero-knowledge proof in which both sides
demonstrate knowledge of the password. Neither side ever transmits anything an offline attacker
could use to mount a dictionary attack against a captured exchange.
\end{definition}

\textbf{WPA3}, required on newly certified Wi-Fi devices since July 1, 2020, closes exactly the gap
WPA2's handshake left open.\srccite{3.2}{3} SAE means capturing the handshake gains an attacker
nothing to brute-force offline. Forward secrecy is built in as well: each session negotiates fresh
keys, so compromising today's session key does not unlock yesterday's captured traffic.
WPA3-Enterprise raises the minimum cipher strength to 192-bit AES, aligned with the NSA's Commercial
National Security Algorithm (CNSA) Suite used in high-security government networks -- the successor
to the older ``Suite B'' name still seen in some vendor documentation.\srccite{3.2}{4}
\textbf{Enhanced Open (OWE)} extends real protection to open networks like a coffee shop's,
negotiating an unauthenticated but still encrypted connection. An eavesdropper in the same room
cannot read your traffic even though there is no password at all.

\begin{lstlisting}[language=Python, caption={Evaluating a Wi-Fi configuration against baseline security criteria}]
def evaluate_wifi_config(cfg):
    """cfg: dict like {"protocol": str, "password_length": int,
       "guest_isolated": bool, "admin_password_default": bool}.
       Returns a sorted list of security gap labels."""
    gaps = []
    if cfg["protocol"] in ("wep", "wpa-tkip"):
        gaps.append("broken-protocol")
    elif cfg["protocol"] == "wpa2" and cfg["password_length"] < 20:
        gaps.append("weak-psk-length")
    if not cfg.get("guest_isolated", False):
        gaps.append("guest-not-isolated")
    if cfg.get("admin_password_default", False):
        gaps.append("default-admin-credentials")
    return sorted(gaps)
\end{lstlisting}

Evaluating a real deployment comes down to a short checklist. Prefer WPA3, accept WPA2-AES as a
floor, and reject TKIP and WEP outright. Use long, random WPA2-PSK passwords (twenty characters or
more) where WPA3 is unavailable, since every added character multiplies the offline search space.
Isolate guest devices on a separate SSID and VLAN that cannot reach internal resources. Keep
firmware current: protocol-level bugs like KRACK, the 2017 key-reinstallation flaw that let an
attacker force reuse of an already-negotiated WPA2 encryption key, are patched at the firmware
layer, not the password layer.\srccite{3.2}{5}

Two further settings matter beyond the protocol choice itself. \textbf{Protected Management Frames
(PMF)}, standardized as 802.11w and mandatory under WPA3, cryptographically protect the
deauthentication and disassociation frames an attacker could otherwise forge to knock a client off
the network at will. Without PMF, an attacker who cannot break the encryption can still force a
client to disconnect and reconnect, creating a fresh handshake to capture. \textbf{WPS (Wi-Fi
Protected Setup)}, a convenience feature meant to simplify pairing a new device, ships enabled by
default on many consumer routers and accepts an eight-digit PIN. A well-documented brute-force flaw,
first published in 2011, recovers that PIN in a matter of hours regardless of how strong the
underlying WPA2 password is.\srccite{3.2}{6} Disabling WPS removes that shortcut entirely.

\begin{misconception}
\textbf{``Hiding the SSID makes the network secure.''} A hidden SSID is still broadcast in probe
responses whenever a client checks for it; any passive scanner finds it within seconds. SSID
hiding is not a security control, just a minor inconvenience.
\end{misconception}

\section{Managerial Controls and Security Policy}

Not every control that protects a network is hardware or software. Security controls are commonly
grouped into three families, and correctly sorting a given control into its family is itself an
AP-level skill. \textbf{Technical (logical) controls} are implemented in hardware or software and
enforce policy automatically -- a firewall rule, WPA3 encryption, an access-control list.
\textbf{Physical controls} restrict physical access, as in Chapters~\ref{ch:2.2}
and~\ref{ch:2.3} -- a locked door, a badge reader.

\begin{definition}{Managerial (administrative) control}
A documented policy, procedure, or process governing human behavior and organizational process,
enforced through training, review, and consequences rather than automatically by hardware or
software.
\end{definition}

A written \textbf{Acceptable Use Policy (AUP)} does not stop an employee from emailing a customer
list to a personal account the way a data-loss-prevention filter might. It sets the expectation
instead, and the organization relies on training, review, and disciplinary consequence to make it
effective. That does not make it a weaker control, only a different kind: an employee who emails
company data to a personal account has violated the AUP even though no technical control was ever
bypassed.

A \textbf{least-privilege policy} establishes that every account should receive only the minimum
access its role requires. As a managerial control, this means the documented process itself -- who
approves an access request, how often access is re-certified, how it is revoked when someone
changes roles. The technical control, by contrast, is the ACL that actually enforces a given
permission. Without the managerial layer, permissions accumulate through \textbf{privilege creep}:
access granted for a since-ended project or a past role that nobody ever revokes.

Four other managerial controls recur across security frameworks. \textbf{Standards and baselines}
translate a policy's intent into a specific, checkable configuration -- a password-complexity
standard, a minimum patch-level baseline -- so ``secure the network'' becomes a requirement an
auditor can actually verify. \textbf{Change control} requires that any modification to a production
system be proposed, reviewed, and approved before it is made, so a well-intentioned configuration
change cannot silently reopen a closed vulnerability. \textbf{Vendor and third-party agreements}
extend the same managerial logic outside the organization's own staff: a supplier or contractor
with network access is bound, contractually, to the same security expectations an employee already
accepted. \textbf{Audit} closes the loop -- an independent review confirming that policy, baseline,
and configuration still agree with each other, catching drift that accumulates silently between one
review and the next.

\textbf{Security awareness training} closes a gap no purely technical control can. An employee
trained to recognize a phishing email functions as a control in exactly the sense a spam filter
does, just implemented through people rather than software. It is the layer most directly connected
back to Chapter~\ref{ch:1.1}'s social-engineering techniques.

\begin{lstlisting}[language=Python, caption={Classifying security scenarios by control family}]
def classify_control(description):
    """Returns "managerial", "technical", or "physical" for a scenario
    description, using keyword signals."""
    lowered = description.lower()
    physical_words = {"lock", "badge", "camera", "guard", "door", "fence"}
    technical_words = {"firewall", "encryption", "acl", "vpn", "antivirus", "wpa"}
    if any(w in lowered for w in physical_words):
        return "physical"
    if any(w in lowered for w in technical_words):
        return "technical"
    return "managerial"
\end{lstlisting}

A policy that exists only on paper, with no training and no consequence for violation, is close to
useless -- employees never learn it exists, or learn that ignoring it costs nothing. A technical
control deployed with no policy explaining its purpose fares little better, generating confused
workarounds from people who do not understand why an access request was denied. NIST's SP 800-53
control catalog reflects this explicitly, placing Awareness and Training (AT) and Planning (PL)
families alongside technical families like System and Communications Protection
(SC).\srccite{3.2}{7} Policy and technology are meant to be read together, not as substitutes for
each other.

\begin{misconception}
\textbf{``A policy is not a real control unless it's enforced by software.''} Managerial controls
are a distinct, legitimate control family. A signed AUP, a documented access-review cadence, or
completed awareness training all count as controls in their own right, even though none of them
are software.
\end{misconception}

\section{Layered Network Defense}

Every control covered so far in this unit -- DNS protections, ARP/MITM defenses, encryption against
packet capture, WPA3, and managerial policy -- addresses one specific threat at one specific point.
None of them alone protects an organization completely, because every individual control has a
known weakness, an undiscovered one, or a misconfiguration waiting to happen.

\begin{definition}{Defense in depth (layered defense)}
The principle that an organization should never rely on a single control to stop an attack.
Multiple independent layers should each be capable of stopping or slowing an attacker, so a failure
at any one layer does not by itself produce a full compromise.
\end{definition}

A commonly taught model orders these layers from the outside in. The \textbf{perimeter layer} --
firewalls, intrusion prevention, DNS filtering -- tries to stop obviously malicious traffic before
it ever reaches internal systems. The \textbf{network layer} limits lateral movement once something
gets past the perimeter: VLANs, internal firewalls between segments, and WPA3 protecting the
wireless portion specifically. The \textbf{host layer} assumes the network layer may fail and
protects individual devices directly -- endpoint antivirus, host firewalls, OS patching. The
\textbf{application layer} lives inside the software itself: input validation, secure coding
practice, a web application firewall. The \textbf{data layer} is innermost and protects the data
regardless of who has network access at all -- encryption at rest and in transit, data-loss
prevention tooling. Even total failure of every outer layer still leaves an attacker holding data
they cannot use. And the \textbf{human/managerial layer} does not occupy a network location at all;
it spans every other layer. Awareness training, the AUP, and least-privilege review are what keep
every technical layer correctly configured, and they are frequently what catches an attack that
slipped past every automated control -- an employee reporting a phishing email the spam filter
missed, for example.

\begin{figure}[h]
\centering
\begin{tikzpicture}[
  ring/.style={draw=kickstartermaroon, thick},
]
\foreach \r/\clr/\y in {3.0/kickstartercream/0, 2.4/white/0, 1.8/kickstartermaroon!10/0, 1.2/kickstartermaroon!20/0, 0.6/kickstartermaroon!35/0}
  \draw[ring, fill=\clr] (0,0) circle (\r);
\node[font=\tiny, align=center] at (0,0) {Data};
\node[font=\tiny] at (0,0.9) {Application};
\node[font=\tiny] at (0,1.5) {Host};
\node[font=\tiny] at (0,2.1) {Network};
\node[font=\tiny] at (0,2.7) {Perimeter};
\draw[kickstartermaroon, thick, dashed] (-4.3,-3.4) rectangle (4.3,3.4);
\node[font=\scriptsize\sffamily, kickstartermaroon, anchor=south west] at (-4.3,-3.4) {Human / managerial layer -- spans every ring above};
\end{tikzpicture}
\caption{Defense in depth as concentric layers protecting the data at the center; the
human/managerial layer is not a ring at all, but spans every one of them.}
\label{fig:3.2-defense-layers}
\end{figure}

\begin{example}[Mapping this unit's controls onto the model]
DNSSEC and source-port randomization sit at the perimeter/network layer, protecting name resolution
before traffic even reaches its destination. Dynamic ARP Inspection and VLAN segmentation against
ARP spoofing sit at the network layer. TLS/HTTPS -- the control that makes captured packets useless
even when capture itself succeeds -- sits at the data layer, protecting the payload regardless of
which layer the capture happened at. WPA3 sits at the network layer, specifically covering the
wireless segment. Managerial controls sit at the human layer, reinforcing all of the above. Every
control this unit has introduced has exactly one home in this picture, which is the entire point of
studying them together rather than in isolation.
\end{example}

\begin{lstlisting}[language=Python, caption={Assigning security controls to their defense-in-depth layer}]
LAYER_KEYWORDS = {
    "perimeter": {"firewall", "ips", "dns filtering"},
    "network": {"vlan", "segmentation", "wpa3", "arp inspection"},
    "host": {"antivirus", "edr", "patch", "host firewall"},
    "application": {"input validation", "waf", "secure coding"},
    "data": {"encryption", "dlp", "data classification"},
    "human": {"training", "aup", "least privilege", "incident response"},
}

def assign_layer(control_description):
    """Returns the defense-in-depth layer name matching keywords in
    control_description, or "unclassified" if none match."""
    lowered = control_description.lower()
    for layer, keywords in LAYER_KEYWORDS.items():
        if any(k in lowered for k in keywords):
            return layer
    return "unclassified"
\end{lstlisting}

A reasonable question is why an organization would not simply deploy the single strongest control
available and skip the rest. The answer is that different controls defend fundamentally different,
non-overlapping attack vectors. WPA3 does nothing against DNS cache poisoning, DNSSEC does nothing
against a local ARP-spoofing MITM, and strong encryption does nothing against an employee who hands
over a password directly to a convincing phone call. An attacker only needs the one vector that is
weakest at any given layer -- which is exactly why depth, not any single control's strength, is what
determines an organization's real posture.

\begin{misconception}
\textbf{``The strongest single control is better than several weaker ones.''} Different controls
defend against different, non-overlapping attack vectors. An attacker only needs to find the one
gap a single-control strategy leaves open; depth closes gaps a single strong control cannot even
see.
\end{misconception}

\section*{Review Questions}
\begin{reviewquestions}
\item Explain why capturing a WPA2 four-way handshake gives an attacker something useful (an
offline dictionary attack target) while capturing a WPA3 SAE handshake does not.
\item An organization's guest Wi-Fi shares an SSID and VLAN with its internal staff network.
Identify which specific configuration gap this represents and explain the risk it creates using
the flat-network reasoning from Chapter~\ref{ch:3.1}.
\item Classify each of the following as managerial, technical, or physical: a badge reader on a
server room door, a signed Acceptable Use Policy, and a firewall rule blocking inbound Telnet.
Justify each classification.
\item Explain why a technical control and the managerial policy governing it are not
interchangeable, using the least-privilege example from this chapter.
\item An attacker defeats a perimeter firewall through a previously unknown vulnerability. Using
the defense-in-depth model, explain which layer should limit the damage next, and why relying on
the perimeter alone would have made this single failure catastrophic instead of merely serious.
\item Explain why the human/managerial layer is drawn spanning every ring in the defense-in-depth
diagram rather than as one more concentric ring of its own, and give one concrete example from
this chapter of a managerial control catching something every technical layer missed.
\item An organization deploys WPA3 network-wide and concludes its wireless traffic is now fully
protected. Using the layered-defense reasoning from this chapter, explain what this conclusion
overlooks about DNS cache poisoning and ARP spoofing, both covered in Chapter~\ref{ch:3.1}.
\item A company patches its servers on an unpredictable, ad-hoc schedule, and no one signs off
before a change goes live. Which managerial control is missing, and what specific risk does its
absence create?
\item A consumer-grade access point is configured for WPA3 with a strong password, but WPS is left
enabled and Protected Management Frames are left off. Explain what each of these two oversights
still leaves open, even though the core wireless protocol itself is not at fault.
\end{reviewquestions}

\begin{answerkey}
\item A captured WPA2 four-way handshake gives an attacker material to run an offline dictionary or
brute-force attack against the pre-shared key. There is no rate limit at all, since the guessing
happens entirely on the attacker's own hardware, so a weak password is recovered this way regardless
of AES's own strength. WPA3's SAE (Dragonfly) handshake is a zero-knowledge proof: both sides
demonstrate knowledge of the password without ever transmitting anything an offline attacker could
use to test guesses against a captured exchange. Capturing the handshake yields nothing to
brute-force.
\answerlabel{3.2.B}
\item This is a flat-network configuration: guest and internal staff traffic share the same
broadcast domain instead of being isolated onto a separate SSID and VLAN. Following the reasoning
from Chapter~\ref{ch:3.1}, a compromised guest device on a shared segment has direct visibility into
every other device sharing it. A single infected visitor laptop can scan and potentially reach
internal staff workstations, printers, and servers with no boundary in the way.
\answerlabel{3.2.A}
\item The badge reader is a physical control -- it restricts physical access to a space using a
tangible mechanism. The signed Acceptable Use Policy is a managerial control -- it governs human
behavior through documentation, training, and consequence rather than through hardware or software
enforcement. The firewall rule is a technical control -- it is implemented in software and enforces
the policy automatically, with no human decision required at the moment traffic arrives.
\answerlabel{3.2.A}
\item The technical control (an ACL entry granting a specific permission) and the managerial policy
governing it (who approves the access request, how often it is re-certified, how it is revoked when
a role changes) answer different questions. The ACL enforces a permission mechanically, while the
policy governs whether that permission should still exist at all. Without the managerial layer,
permissions accumulate as privilege creep -- access granted for a since-ended project or a past role
that nobody with authority ever reviews or revokes. The technical control itself keeps functioning
exactly as configured the whole time.
\answerlabel{3.2.A}
\item In the defense-in-depth model, the network layer -- VLANs and internal firewalls limiting
lateral movement -- should be what stops the attacker next, since the perimeter's failure only grants
entry, not free reach across the internal network. Relying on the perimeter alone would have made
this catastrophic because a flat internal network gives a single successful breach immediate,
unopposed access to every other system. With network-layer segmentation in place, the same breach
is contained to whatever segment the attacker first landed in, turning one failure into a serious
but bounded incident rather than a total compromise.
\answerlabel{3.2.A}
\item The human/managerial layer is not tied to any single network location or technology the way
perimeter, network, host, application, and data controls are. Awareness training and least
privilege review affect how correctly every one of those other layers is configured and operated,
so a ring drawn at only one radius would understate its reach. A concrete example: security
awareness training produces an employee who reports a phishing email the spam filter (a technical,
perimeter-adjacent control) missed entirely, catching an attack no purely technical layer in the
model stopped.
\answerlabel{3.2.A}
\item WPA3 fully addresses the wireless segment specifically -- protecting the radio link between a
device and its access point -- but DNS cache poisoning and ARP spoofing operate at entirely
different points that WPA3 has no relationship to at all. Cache poisoning exploits the resolver's
inability to authenticate a DNS reply, and ARP spoofing exploits the local network's lack of
address-binding authentication; neither is touched by WPA3's link-layer encryption. Concluding that
wireless traffic is "fully protected" mistakes strength at one layer for coverage of every layer.
The organization still needs DNSSEC, Dynamic ARP Inspection, and segmentation to close gaps WPA3
was never designed to close in the first place.
\answerlabel{3.2.B}
\item The company is missing change control. Without a required propose-review-approve step before
a production change goes live, a well-intentioned patch can reintroduce a vulnerability that was
already closed, or collide with another change already in progress, with no independent check
catching it before it reaches production.
\answerlabel{3.2.A}
\item WPS leaves an eight-digit PIN exposed to a well-documented brute-force flaw that recovers it
in hours, handing an attacker network access through a path that has nothing to do with the WPA3
password's own strength -- the strong password never gets tested at all. The missing Protected
Management Frames leave deauthentication and disassociation frames unauthenticated, so an attacker
can still forcibly disconnect a legitimate client at will, whether to deny service or to force a
fresh handshake to capture. Both gaps sit outside the core WPA3 protocol, which is exactly why
configuring a strong protocol is necessary but not sufficient on its own.
\answerlabel{3.2.B}
\end{answerkey}

\section*{Sources}
\begin{enumerate}[label=\lbrack\arabic*\rbrack, leftmargin=2.2em]
\item \srcitem{3.2}{1}Scott Fluhrer, Itsik Mantin, and Adi Shamir, ``Weaknesses in the Key Scheduling
Algorithm of RC4,'' \emph{Selected Areas in Cryptography (SAC 2001)}, Springer LNCS 2259, 2001 --
the original paper demonstrating the flaw in WEP's RC4 key scheduling.
\item \srcitem{3.2}{2}Erik Tews, Ralf-Philipp Weinmann, and Andrei Pyshkin, ``Breaking 104 Bit WEP
in Less Than 60 Seconds,'' IACR Cryptology ePrint Archive, Report 2007/120 -- the practical attack
recovering a WEP key from as few as 40{,}000 captured frames (50\% success probability; 85{,}000
frames for 95\%).
\item \srcitem{3.2}{3}Wi-Fi Alliance, announcement that WPA3 certification became mandatory for all
newly Wi-Fi CERTIFIED devices effective July 1, 2020.
\item \srcitem{3.2}{4}National Security Agency, Commercial National Security Algorithm (CNSA) Suite
specification, the successor to the earlier ``Suite B'' cryptographic requirements referenced by
WPA3-Enterprise's 192-bit security mode.
\item \srcitem{3.2}{5}Mathy Vanhoef and Frank Piessens, ``Key Reinstallation Attacks: Forcing Nonce
Reuse in WPA2,'' \emph{Proceedings of the 2017 ACM SIGSAC Conference on Computer and Communications
Security (CCS)}, October 2017 -- the paper disclosing the KRACK vulnerability.
\item \srcitem{3.2}{6}Stefan Viehb\"ock, ``Wi-Fi Protected Setup PIN Brute Force Vulnerability,''
published December 27, 2011 (also tracked as CERT/CC vulnerability note VU\#723755) -- the
disclosure of the WPS PIN brute-force flaw.
\item \srcitem{3.2}{7}National Institute of Standards and Technology, Special Publication 800-53
Revision 5, ``Security and Privacy Controls for Information Systems and Organizations,'' September
2020.
\end{enumerate}
